\documentclass[12pt]{article}

% --------------------
% PAQUETES BÁSICOS
% --------------------
\usepackage[spanish,es-noshorthands]{babel}
\usepackage[utf8]{inputenc}      % Si usas pdflatex
\usepackage[T1]{fontenc}
\usepackage{lmodern}
%\usepackage[authoryear]{natbib}
\usepackage{comment}
\usepackage{amsmath, amssymb, amsthm}
\usepackage{geometry}
\usepackage{booktabs}   
\usepackage{setspace}
\usepackage{graphicx}
\usepackage{tikz}
\usepackage{tikz-cd}
\usepackage[
  colorlinks=true,
  linkcolor=blue,
  citecolor=blue,
  urlcolor=red
]{hyperref}

\newtheorem{thm}{Teorema}[section]
\newtheorem{dfn}[thm]{Definición}
\newtheorem{lem}[thm]{Lema}
\newtheorem{ej}[thm]{Ejemplo}
\newtheorem{cor}[thm]{Corolario}
\newtheorem{obs}[thm]{Observación}
\newtheorem{con}[thm]{Conjetura}
\newtheorem{prop}[thm]{Proposición}

\DeclareMathOperator{\Frm}{Frm}
\DeclareMathOperator{\Fil}{Fil}
\DeclareMathOperator{\Ord}{Ord}
\DeclareMathOperator{\Top}{Top}
\DeclareMathOperator{\Sp}{sp}
\DeclareMathOperator{\pt}{pt}
\DeclareMathOperator{\id}{id}
\DeclareMathOperator{\tp}{tp}
\DeclareMathOperator{\fd}{fd}
% --------------------
% CONFIGURACIÓN
% --------------------
\geometry{margin=2.5cm}
\onehalfspacing

% --------------------
% DATOS DEL DOCUMENTO
% --------------------
\title{Reporte mensual de actividades}
\date{1-septiembre-2026 a 30-septiembre-2026}
\author{Juan Carlos Monter Cortés}

% --------------------
% INICIO DEL DOCUMENTO
% --------------------
\begin{document}

\maketitle
% \tableofcontents

\section{Actividades realizadas}

Siguiendo lo presentado en \cite{sexton2003point}, para $A$ un marco, $F\in A^\wedge$ y $\alpha\in \Ord$ decimos que:
\begin{itemize}
    \item $F$ es $\alpha$-eficiente ($\alpha$-neat), si
    \[
    x\vee d=1,\quad \text{para toda }x\in F,
    \]
    donde 
    \[
    d=f_F^\alpha(0)\quad\text{y}\quad f_F=\dot{\bigvee}\{v_a\mid a\in F\}.
    \]
    \item $A$ es $\alpha$-eficiente si para todo $F\in A^\wedge$, $F$ es $\alpha$-eficiente.
    \item $A$ es eficiente si es $\alpha$-eficiente para algún $\alpha$.
\end{itemize}

Adicionalmente, se puede verificar que $A$ es eficiente si se cumple que
\[
f_F^\infty=v_F=u_d, \quad\text{para todo }F\in A^\wedge \text{ y }d=v_F(0).
\]

También, sabemos que un marco $A$ es KC si para todo $F\in A^\wedge$ y $j\in NA$ tal que 
\[
j\in [v_F, w_F], \text{ existe } a\in A\text{ tal que } j=u_a.
\]

De esta manera, si $A$ es un marco KC y $F\in A^\wedge$, entonces $v_F=u_a$. Evaluando en 0, 
\[
v_F(0)=u_a(0)=a\quad\text{ y }\quad a=d.
\]

Por lo tanto, KC implica Eficiente.

Tomando $S=\pt A$, tenemos la siguiente caracterización:
\[
A \text{ es eficiente }\iff S \text{ es empaquetado y apilado}.
\]
donde apilado es la noción espacial presentada en \cite{sexton2006ordinal}.

Así, considerando lo presentado en reportes anteriores, podemos preguntarnos cuál es la relación de los marcos KC y los marcos eficientes con los marcos 
apilados y fuertemente apilados. Lo anterior, es parte de los objetivos del presente reporte. Además, presentamos cómo la propiedad KC se relaciona con 
la propiedad $\mathbf{(sfit)}$ y algunas de las consecuencias de la interacción entre ambas.

\begin{prop}\label{KCSS}
    Sean $A$ un marco y $(F,Q)$ un par HM. Si $A$ es KC, entonces las nociones de apilado y fuertemente apilado son equivalentes.
\end{prop}

\begin{proof}
    Ya que KC implica $T_1$, para cualquier par HM $(F,Q)$ se satisface que 
    \[
    \mathcal{U}_*\circ [Q']\circ \mathcal{U}=w_F.
    \]

    Además, por la definición de KC, el intervalo de admisibilidad es de la forma
    \[
    [v_F=u_d, w_F=u_a],
    \]
    donde $d=v_F(0)$ y $a=w_F(0)$.

    Por lo tanto, debemos probar que $u_d=u_a$. Así,
    \begin{equation}\label{ud=ua}
    \text{Fuertemente apilado }\iff u_d=u_a\iff a=d\iff v_F(0)=w_F(0)\iff \text{ Apilado}.
    \end{equation}
\end{proof}

Recordemos que un espacio $S$ es empaquetado si para todo $Q\in\mathcal{Q}S$, $Q\in \mathcal{C}S$, donde $\mathcal{Q}S$ y $\mathcal{C}S$ son 
las familias de subconjuntos compactos saturados y subconjuntos cerrados, respectivamente.

Dado que KC implica eficiente y eficiente implica empaquetado, entonces 
\[
\text{KC}\Rightarrow \text{Empaquetado},
\]
es decir, para $S$ un espacio topológico, si $\mathcal{O}S$ es KC, el espacio $S$ es empaquetado.

El reciproco de lo anterior no es cierto. De hecho, si consideramos un espacio empaquetado $S$, tenemos
\[
Q\in \mathcal{Q}S\Rightarrow Q'\in \mathcal{O}S.
\]

Luego, considerando $F\in \mathcal{O}S^\wedge$, el filtro abierto asociado a $Q$, se cumple que
\[
v_F\leq [Q']=u_{Q'}.
\]
y como empaquetado implica $T_1$, entonces $[Q']=w_F$.\\

Por lo tanto, si $S$ es fuertemente apilado, $v_F=u_{Q'}=w_F$, es decir, 
\[
\mathcal{O}S \text{ es KC }\iff S\text{ es empaquetado y fuertemente apilado}.
\]

\begin{prop}
    Si $A$ es un marco espacial, entonces las nociones KC y Eficiente son equivalentes.
\end{prop}

\begin{proof}
    En reportes anteriores probamos que si $\mathcal{O}S$ es eficiente, entonces $S$ es fuertemente apilado. Además, 
    por la Proposición \ref{KCSS},
    \[
    \begin{split}
    \mathcal{O}S \text{ es KC }&\iff S\text{ es empaquetado y fuertemente apilado}\\
    &\iff \mathcal{O}S\text{ es eficiente}.
    \end{split}
    \]
\end{proof}

El resultado anterior nos dice que, en el caso espacial, KC y Eficiente son en escencia la misma propiedad. Ahora, debemos
preguntarnos si ocurre el mismo fenomeno para un marco arbitrario. Para poder responder esto, primero observamos la relación
de KC y F. apilado.

Recordemos que un marco $A$ es subajustado ($(\mathbf{sfit})$) si para cualesquiera $a,b\in A$, con $a\nleq b$, se cumple 
\[
\exists\, c\in A\text{ tal que }a\vee c=1\neq b\vee c.
\]

En \cite{simmons2006regularity,picado2021separation} caracterizan esta clase de marcos. Haciendo uso de ellos, 
podemos afirmar lo siguiente.

\begin{prop}\label{KCsfit}
    Si $A$ es un marco KC y $\mathbf{(sfit)}$, entonces $A$ es fuertemente apilado.
\end{prop}

\begin{proof}
    Dado que en los marcos KC todo cociente compacto es cerrado, debemos verificar que $u_d=u_a$. Por \eqref{ud=ua}, lo anterior 
    es equivalente a verificar que $d=a$. Por hipótesis, $A$ es $\mathbf{(sfit)}$ y, por \cite[Teorema 3.5]{simmons2006regularity}, 
    se cumple quue $d=a$. Por lo tanto, $A$ es fuertemente apilado.
\end{proof}

En \cite[Proposición 2.5.3]{picado2021separation} mencionan que si un marco es $T_1$ y $\mathbf{(sfit)}$, entonces este es espacial. Este hecho
y el ejemplo presentado en el reporte anterior de un marco $T_1$ y fuertemente apilado que no es espacial, nos permite afirmar lo siguiente.

\begin{prop}
    Si un marco es KC, entonces no es $\mathbf{(sfit)}$
\end{prop}

\begin{proof}
    A manera de contradicción, supongamos que KC implica $\mathbf{(sfit)}$. Así, por la Proposición \ref{KCsfit}, tenemos que $KC$ siempre implica fuertemente 
    apilado. Luego, como fuertemente apilado implica apilado,
    \[
    KC\Rightarrow \text{ Apilado}+T_1\Rightarrow \text{ Espacial},
    \]
    lo cual es una contradicción, pues existen marcos $T_1$ y apilados que no son espaciales. 
\end{proof}

Observemos que si KC implica fuertemente apilado, se cumple que 
\[
[v_F, w_F]=\{u_d\}.
\]

Entonces, similarmente a como ocurre en el caso espacial, las nociones de KC y Eficiente son equivalentes para marcos abitrarios. Por lo tanto, dentro de las actividades 
futuras se encuentra el verificar si existen marcos KC que no son fuertemente apilados.

\bibliographystyle{amsalpha}
\bibliography{research2}

\end{document}